\documentclass[11pt]{article}
\usepackage[margin=1in]{geometry}
\usepackage{amsmath,amssymb}
\usepackage{booktabs}
\usepackage[dvipsnames]{xcolor}
\usepackage{hyperref}
\hypersetup{colorlinks=true,linkcolor=blue!50!black,urlcolor=blue!50!black,citecolor=blue!50!black}

\title{Do Fairness Mitigations Beat Test-Set Noise?\\
Person-Level Uncertainty for a Preregistered\\
Policy Comparison on UCI Adult}
\author{David Lindgreen}
\date{October 9, 2026 \\[2pt] \small Version 0.1.0 \quad Preprint, not peer reviewed}

\emergencystretch=3em
\begin{document}
\maketitle

\begin{abstract}
Fairshift Lab compares four decision policies (a fixed threshold, a cost-selected threshold, reweighed training and group-specific
thresholds) on the UCI Adult census-income benchmark~\cite{becker1996}, and reported that mitigation lowered the observed
equalized-odds gap, for example from 0.037 to 0.017. Its uncertainty, however, covers only five random 80/20 splits of the
development file; the held-out test file (15{,}060 people in the primary cohort) is never resampled, so the report cannot say
whether a gap difference of this size is larger than the sampling noise of the test people themselves. We re-ran all four frozen
registries (they reproduce exactly) and added a paired person-level bootstrap and a permutation noise floor. The equalized-odds
gaps are imprecise (the baseline's 95\% interval is 0.010 to 0.080) and sit at about twice a no-association floor of 0.014; in none of
60 paired comparisons across four cohort definitions and five split seeds does the interval for the baseline-minus-policy gap
difference exclude zero. Only the accuracy cost of group thresholds is distinguishable from noise (18 of 20 cells). The preregistered
conclusion that no policy is uniformly best is unaffected; statements that a mitigation ``reduced'' the gap are not supported at this
test-set size.
\end{abstract}

\section{Question}
Group-fairness measurements are estimates from a finite sample of people, and a difference of two such estimates has its own
sampling error~\cite{besse2022}. Fairshift Lab's external study fixed its design before interpreting held-out results and reported the
mean and the 2.5th--97.5th range of each metric across five development splits. That range describes how much the answer depends on
which people were used for training and tuning. It does not describe how much it depends on which people were in the test file. Here
we measure that second source of uncertainty.

\section{Frozen study and reproduction}
The study uses the provider-coded binary sex field, the five numeric features and the income label (above USD 50{,}000 in 1994, which is
not a measure of merit or need). Policies are evaluated at false-negative cost 1 against false-positive cost 1; thresholds are chosen on a
tuning split only. We re-ran the synthetic study, the policy study, the robustness study and the external study from the pinned
sources: all four registries reproduce the committed files with no differences, and the new analysis recovers the frozen
equalized-odds means exactly for every policy.

\section{Post-hoc audit}
For each of four cohort/missingness definitions and each of five split seeds we re-created the four policies' held-out decisions and drew
2{,}000 bootstrap resamples of the test people, with the same resample applied to all policies so that differences are paired. We also
computed a \emph{noise floor}: the mean equalized-odds difference after randomly permuting the protected attribute among test people
500 times. Because the statistic is the larger of two absolute rate gaps, it is positive even when the attribute carries no information,
and the floor shows how large that bias is. This analysis was not preregistered and all of it is reported.
\begin{table}[h]
\centering\small
\begin{tabular}{lccc}
\toprule
Policy (primary cohort) & Equalized-odds gap & 95\% interval & Noise floor\\
\midrule
Baseline threshold 0.5 & 0.037 & 0.010 to 0.080 & 0.014\\
Cost-selected threshold & 0.039 & 0.011 to 0.082 & 0.014\\
Reweighed training & 0.021 & 0.004 to 0.063 & 0.014\\
Group thresholds & 0.017 & 0.004 to 0.059 & 0.014\\
\bottomrule
\end{tabular}
\caption{Complete-case provider-eligible cohort (15{,}060 test people), means over five split seeds.}
\end{table}
\begin{table}[h]
\centering\small
\begin{tabular}{lcc}
\toprule
Baseline minus policy & Accuracy & Equalized-odds gap\\
\midrule
Reweighed training & $-0.001$ ($-0.003$ to $+0.001$) & $+0.016$ ($-0.014$ to $+0.031$)\\
Group thresholds & $+0.004$ ($+0.002$ to $+0.007$) & $+0.020$ ($-0.040$ to $+0.055$)\\
\bottomrule
\end{tabular}
\caption{Paired differences with 95\% person-level bootstrap intervals (primary cohort, mean over seeds).}
\end{table}
Across all four cohorts and five seeds (60 paired comparisons of a mitigation with the baseline) the equalized-odds interval never
excludes zero. Point estimates range from 0.69 to 3.4 times the noise floor, averaging about 2.0. The accuracy cost of group thresholds,
by contrast, is small but resolved: the interval excludes zero in 18 of 20 cells.

\section{Discussion}
The registered point estimates are correct and unchanged; what changes is how strongly they can be read. With about 15{,}000 test
people, and far fewer in the smaller group's positive class, a gap of 0.02 versus 0.04 cannot be told apart from
resampling noise. The frozen hypotheses (no policy uniformly best; conclusions sensitive to cohort and missingness) concern rankings that
vary, so they stand; the prose that mitigation ``reduced the observed gap'' should be read as describing this one test file.
Fairness-mitigation comparisons on benchmarks of this size need person-level intervals, and any claim of reduction should be accompanied
by the noise floor for the statistic used.

\section{Limitations}
The bootstrap treats the held-out people as the sampling unit and ignores training-set variation beyond the five split seeds; it describes
an observational reference table, not a deployment population. Sex is a provider-coded binary field. The audit uses one cost setting
(false-negative cost 1); other costs are in the frozen registry but were not resampled. Intervals are percentile bootstrap intervals
without multiplicity adjustment.

\section*{Reproducibility and AI assistance}
Code, frozen registries, the audit and tests:\\
\url{https://github.com/lindgreendavid/fairshift-lab}\\
The analysis code, tests and drafts of this text were produced with substantial assistance from an AI system (Claude, Anthropic);
the author is responsible for the content.

\begin{thebibliography}{9}
\bibitem{becker1996} B.~Becker and R.~Kohavi, Adult, UCI Machine Learning Repository (1996), doi:10.24432/C5XW20.
\bibitem{besse2022} P.~Besse, E.~del Barrio, P.~Gordaliza, J.-M. Loubes, L.~Risser, A survey of bias in machine learning through
the prism of statistical parity, \emph{Am. Stat.} \textbf{76}, 188--198 (2022).
\end{thebibliography}
\end{document}
